\subsection{Warped product cones and flat completion}
\label{sec:warped-gauduchon-completion}

The invariant cones above have complete unit-volume hypersurfaces. On the product of a compact K\"ahler manifold with a curve, a nonconstant coefficient in the fiber direction gives cones of nonclosed backgrounds whose unit-volume hypersurfaces are flat and may instead have a Euclidean strip or half-plane as their metric completion. The volume is the reciprocal of an integral of the reciprocal fiber coefficient, and the variance of that reciprocal coefficient measures the distance to the boundary. The probability measure in this variance is a trace of a minimizing balanced metric.

Let $(B,\kappa)$ be a connected compact K\"ahler manifold of dimension $m\geq1$, with $\int_B\kappa^m=1$. Let $C$ be a connected compact complex curve with a smooth positive $(1,1)$-form $w$ of integral one. Fix a smooth strictly positive probability form $\sigma$ on $B$ and a smooth nonconstant real function $a$ on $B$ with $\min a=0$ and $\max a=M>0$. On $X=B\times C$, of dimension $n=m+1$, put $\mu=\sigma\wedge w$; we omit pullbacks from the notation. For real parameters $b,c,e$, the backgrounds in the real span of $\kappa,aw,w$ are
\[
 \theta_{b,c,e}=b\kappa+Lw,\qquad L=ca+e.
\]

\begin{theorem}\label{thm:warped-gauduchon-volume}\index{Gauduchon variational volume!warped product cones}
This span contains no nonzero smooth $\ddc$ direction, and its intersection with $\mathscr P$ is the cone
\[
 \mathscr C=\{b>0,\ e>0,\ e+cM>0\}.
\]
For $\theta=\theta_{b,c,e}\in\mathscr C$, let $H(c,e)=\int_B\sigma/L$. Then
\begin{equation}\label{eq:warped-gauduchon-volume}
 \mathcal V_\mu(\theta)=\frac{nb^m}{H(c,e)}.
\end{equation}
The minimizing Gauduchon metrics are the constant multiples of the balanced metric
\[
 \gamma=L^{-1/m}(A+Lw),\qquad
 A=b\kappa+\ddc\phi>0,\qquad
 A^m=\frac{b^m}{H}\frac{\sigma}{L},
\]
where $\phi$ is a smooth real function on $B$. Backgrounds with $c\ne0$ are nonclosed.
\end{theorem}
\begin{proof}
Pairing a smooth positive representative $\theta+\ddc u$ with the closed semipositive form $\kappa^{m-1}\wedge w$ gives $b>0$, and its integral over the compact fiber $\{y\}\times C$ is $L(y)>0$, by Stokes' theorem. These are exactly the stated inequalities; conversely, they make $\theta$ itself positive. The same tests applied to a smooth $\ddc$ direction give $b=0$ and $ca+e=0$ everywhere, and since $a$ is nonconstant, $c=e=0$. Also $\mathrm{d}\theta=c\,\mathrm{d}a\wedge w$, which is not identically zero when $c\ne0$.

Next we apply \cref{thm:tw-smooth} on $B$ to the closed positive background $b\kappa$ and the density $(b^m/H)\sigma/L$. Since this density has integral $b^m=\int_B(b\kappa)^m$, Stokes' theorem forces the constant in the equation to be one, and we obtain $A$. Put $\Omega=A+Lw$. Since $A^{m+1}=0$ and $(Lw)^2=0$, the binomial expansion gives
\begin{equation}\label{eq:warped-balanced-cofactor}
 \Omega^n=\frac{nb^m}{H}\mu,\qquad
 \gamma^m=L^{-1}A^m+mA^{m-1}\wedge w.
\end{equation}
The first summand is the pullback of a top-degree form on $B$, and the second is closed. Hence $\gamma$ is balanced, in particular Gauduchon, and it is conformal to the solution form $\Omega=\theta+\ddc\phi$ of $\theta$. Since $n\geq2$, \cref{thm:gaud-volume-dual-formula} proves \eqref{eq:warped-gauduchon-volume} and the description of the minimizers.
\end{proof}

The minimizing metric gives the probability identity
\begin{equation}\label{eq:warped-minimizer-trace}
 \frac{Lw\wedge\gamma^m}{b^m}
       =\frac{\sigma\wedge w}{HL}.
\end{equation}
Its marginal $\sigma/(HL)$ on $B$ is the probability measure that governs the Hessian metric. To compute that metric we use the ratio $v=c/e$ and set
\begin{align*}
 v&=c/e\in(-1/M,\infty),
 &Q(v)&=\int_B\frac{\sigma}{1+va},\\
 \mathrm{d}\nu_v&=\frac{\sigma}{Q(v)(1+va)},
 &\ell_v&=\frac{a}{1+va}.
\end{align*}

\begin{proposition}\label{prop:warped-flat-unit-metric}\index{Hessian metric!flat warped cones}
On the unit-volume hypersurface $\mathscr S=\{\theta\in\mathscr C:\mathcal V_\mu(\theta)=1\}$, the functions $U=\log b$ and $v$ are global coordinates, in which
\begin{equation}\label{eq:warped-flat-unit-metric}
 \begin{aligned}
 g_{\mathscr S}&=m\,\mathrm{d}U^2+\rho(v)\,\mathrm{d}v^2,\\
 \rho(v)&=\frac{Q''/Q-2(Q'/Q)^2}{n}
       =\frac2n\operatorname{Var}_{\nu_v}(\ell_v)>0.
 \end{aligned}
\end{equation}
For any $v_*\in(-1/M,\infty)$, the coordinates
\[
 Y=\sqrt m\,U,\qquad
 R(v)=\int_{v_*}^v\sqrt{\rho(t)}\,\mathrm{d}t
\]
define a global isometry onto $\mathbb R\times I$, where $I$ is an open interval with the Euclidean metric. The intrinsic distance is the Euclidean distance in $(Y,R)$.
\end{proposition}
\begin{proof}
Consider an affine variation $\delta=(\delta b,\delta c,\delta e)$ of the coefficients, write $F=\delta L/L$, and note that $\sigma/(HL)=\mathrm{d}\nu_v$ because $HL=Q(v)(1+va)$. Differentiating under the integral sign gives
\[
 D^2\log H(\delta,\delta)
   =2\mathbb E_{\nu_v} F^2-(\mathbb E_{\nu_v} F)^2,
\]
where $\mathbb E_{\nu_v}$ denotes expectation with respect to $\nu_v$. This is positive unless $\delta L=0$ everywhere, which forces $\delta c=\delta e=0$. Since $\log\lambda_\mu=(\log n+m\log b-\log H)/n$,
\[
 g=-D^2\log\lambda_\mu
   =\frac mn\left(\frac{\mathrm{d}b}{b}\right)^2
                  +\frac1nD^2\log H
\]
is positive definite. On the unit-volume hypersurface, the relation $H(c,e)=Q(v)/e=nb^m$ gives
\[
 e=\frac{Q(v)}{n\exp(mU)},\qquad c=ve.
\]
To pull back the Hessian, write $t=\mathrm{d}e/e$ and $z=Q'/Q$. Along an affine path of coefficients one has $v''=-2tv'$, so differentiating $\log H=-\log e+\log Q(v)$ twice gives
\[
 D^2\log H=t^2-2zt\,\mathrm{d}v
                   +(Q''/Q-z^2)\,\mathrm{d}v^2.
\]
On $\mathscr S$ we have $t=z\,\mathrm{d}v-m\,\mathrm{d}U$; substituting this and using $n=m+1$ gives \eqref{eq:warped-flat-unit-metric}. Since $Q'/Q=-\mathbb E_{\nu_v}\ell_v$ and $Q''/Q=2\mathbb E_{\nu_v}\ell_v^2$, the variance formula follows, and nonconstancy of $a$ makes this variance positive. In the coordinates $(Y,R)$ the metric is Euclidean on the convex set $\mathbb R\times I$, so straight segments prove the distance assertion.
\end{proof}

\subsubsection{Lengths of the two ends}

The interval $I$ has two ends, corresponding to $v\to\infty$ and $v\to-1/M$. Whether they lie at finite distance depends on how $a$ attains its extreme values. Assume that each of the extreme sets $\{a=0\}$ and $\{a=M\}$ is a finite disjoint union of smooth compact real submanifolds $Z_i$ of positive codimension, and that near each component, respectively,
\[
 a\asymp\operatorname{dist}(Z_i)^{k_i},\qquad
 M-a\asymp\operatorname{dist}(Z_i)^{k_i},\qquad k_i>0,
\]
where $\asymp$ means comparability up to positive constant factors and $\operatorname{dist}(Z_i)$ is the distance to $Z_i$ for a fixed Riemannian metric on $B$. The function $a$, respectively $M-a$, is then bounded below by a positive constant outside small disjoint tubes around these components. Define, with minima over the components $Z_i$ of the respective extreme set,
\[
 r_0=\min_{\{a=0\}}
       \frac{\operatorname{codim}_{\mathbb R}Z_i}{k_i},\qquad
 r_M=\min_{\{a=M\}}
       \frac{\operatorname{codim}_{\mathbb R}Z_i}{k_i}.
\]

\Needspace{15\baselineskip}
\begin{lemma}\label{lem:warped-endpoint-length}\index{Gauduchon metric!trace variance}
The end $v\to\infty$ of $I$ has finite length exactly when $r_0>1$, and the end $v\to-1/M$ exactly when $r_M>1$. More precisely, at the end $v\to\infty$ put $f=a$, $\epsilon=1/v$ and $r=r_0$; at the end $v\to-1/M$ put $f=M-a$, $\epsilon=-1/v-M$ and $r=r_M$. Let
\[
 H_f=\int_B\frac{\sigma}{f+\epsilon},\quad
 J_f=\int_B\frac{\sigma}{(f+\epsilon)^2},\quad
 K_f=\int_B\frac{\sigma}{(f+\epsilon)^3}.
\]
Then $\epsilon\downarrow0$ at the end, and
\begin{equation}\label{eq:warped-endpoint-variance}
 \rho(v)\,\mathrm{d}v^2=G_f(\epsilon)\,\mathrm{d}\epsilon^2,
 \quad
 G_f=\frac2n\left[\frac{K_f}{H_f}
                    -\left(\frac{J_f}{H_f}\right)^2\right].
\end{equation}
Moreover, $G_f\asymp\epsilon^{-2}$ for $r<1$ and $G_f\asymp\epsilon^{-2}/\log(1/\epsilon)$ for $r=1$, while for $r>1$ the function $\sqrt{G_f}$ is integrable near $\epsilon=0$.
\end{lemma}
\begin{proof}
At the end $v\to\infty$, where $\epsilon=1/v$, we have
\[
 \ell_v=\epsilon-\frac{\epsilon^2}{a+\epsilon},\qquad
 \mathrm{d}v=-\epsilon^{-2}\mathrm{d}\epsilon.
\]
At the end $v\to-1/M$, where $\epsilon=-1/v-M$, we obtain
\[
 \ell_v=\frac{(M+\epsilon)^2}{M-a+\epsilon}-(M+\epsilon),
 \qquad \mathrm{d}v=(M+\epsilon)^{-2}\mathrm{d}\epsilon.
\]
In both cases $\ell_v$ is an affine function of $1/(f+\epsilon)$ with coefficients depending only on $\epsilon$, and \eqref{eq:warped-endpoint-variance} expresses $\rho(v)\,\mathrm{d}v^2$ through the variance of $1/(f+\epsilon)$ under the probability $\sigma/[H_f(f+\epsilon)]$. This probability is $\nu_v$, the marginal in \eqref{eq:warped-minimizer-trace}, because $L$ is a positive multiple of $f+\epsilon$ and positive scaling of $L$ does not change that marginal.

Integration over tubes around the components $Z_i$, splitting the normal radius at $\epsilon^{1/k_i}$, gives
\[
 H_f\asymp
 \begin{cases}\epsilon^{r-1}&r<1,\\
 \log(1/\epsilon)&r=1,\\1&r>1,
 \end{cases}
 \qquad
 K_f\asymp
 \begin{cases}\epsilon^{r-3}&r<3,\\
 \log(1/\epsilon)&r=3,\\1&r>3.
 \end{cases}
\]
The analogous threshold for $J_f$ is two. For a finite sum of tube integrals, the smallest ratio governs.

If $r<1$, choose a component attaining $r$, and two normal annuli at fixed multiples of $\epsilon^{1/k_i}$, so far apart that the two-sided bounds separate their ranges of $f/\epsilon$. Each annulus has probability bounded below, since its weighted mass and $H_f$ both have order $\epsilon^{r-1}$, and the values of $1/(f+\epsilon)$ on the two annuli differ by at least $c\epsilon^{-1}$. The identity
\[
 \operatorname{Var}_\nu T
   =\frac12\iint(T(x)-T(y))^2\,\mathrm{d}\nu(x)\mathrm{d}\nu(y)
\]
applied to $T=1/(f+\epsilon)$ gives the lower bound, and $T\leq\epsilon^{-1}$ gives the upper bound.

If $r=1$, the near annulus has probability at least $c/\log(1/\epsilon)$. In the same tube, the normal radii between $\epsilon^{1/(2k_i)}$ and a fixed small radius carry probability bounded below, by the logarithmic integral. There $1/(f+\epsilon)\leq C\epsilon^{-1/2}$, while $1/(f+\epsilon)\geq c\epsilon^{-1}$ on the near annulus. The variance identity and the bound $G_f\leq2K_f/(nH_f)$ prove the claimed logarithmic estimate. In both cases $r<1$ and $r=1$, the integral of $\sqrt{G_f}$ diverges at $\epsilon=0$.

For $r>1$, $H_f$ tends to a finite positive integral and $G_f\leq C K_f$. For $1<r<3$, $\sqrt{G_f}$ is therefore bounded by $C\epsilon^{(r-3)/2}$, whose exponent exceeds $-1$. At $r=3$ the bound is $C\sqrt{\log(1/\epsilon)}$, and for $r>3$ it is a constant. All these bounds are integrable.
\end{proof}

\begin{theorem}\label{thm:warped-unit-volume-completion}\index{Metric completion!warped unit-volume hypersurface}
Under the hypotheses of \cref{lem:warped-endpoint-length} on the extreme sets, the metric completion of $\mathscr S$ is a closed Euclidean strip if $r_0,r_M>1$, a closed Euclidean half-plane if exactly one of them exceeds one, and the whole Euclidean plane otherwise. Each boundary line of the completion is parametrized by $b>0$, and its points are represented by the forms
\begin{align*}
 \theta_+(b)&=b\kappa+\frac{H_a}{nb^m}a w,
 &H_a&=\int_B\frac{\sigma}{a}\quad(r_0>1),\\
 \theta_-(b)&=b\kappa+\frac{H_-}{nb^m}(1-a/M)w,
 &H_-&=\int_B\frac{\sigma}{1-a/M}\quad(r_M>1).
\end{align*}
These forms are smooth, semipositive, nef, nonclosed and not big, and their variational volume is one. The completion has no other boundary points.
\end{theorem}
\begin{proof}
The completion of the convex product $\mathbb R\times I$ of \cref{prop:warped-flat-unit-metric} is its Euclidean closure. Along a Cauchy sequence, $Y$ converges, so $b$ has a finite positive limit, and $R$ converges to a point of $I$ or to a finite endpoint of $I$. \Cref{lem:warped-endpoint-length} decides which endpoints are finite and gives the three alternatives.

On $\mathscr S$ we have $e=Q(v)/(nb^m)$ and $c=ve$. At the end $v\to\infty$, $vQ(v)\to H_a$ by monotone convergence, which gives the coefficient limit $\theta_+(b)$. At the end $v\to-1/M$, write $1+va=1-a/M+(v+1/M)a$; dominated convergence gives $Q(v)\to H_-$, and hence the limit $\theta_-(b)$. These are the coefficient limits at all finite ends. The forms are semipositive, hence nef, and they are nonclosed. Their fiber coefficients $L_\pm$ vanish somewhere, so they are not big by \cref{prop:warped-current-cones} below. By the definitions of $H_a$ and $H_-$, $\int_B\sigma/L_\pm=nb^m$. For $\epsilon>0$, the form $\theta_\pm(b)+\epsilon(\kappa+w)$ lies in $\mathscr C$, and by \eqref{eq:warped-gauduchon-volume} and monotone convergence its volume
\[
 \frac{n(b+\epsilon)^m}{\int_B\sigma/(L_\pm+\epsilon)}
\]
tends to one as $\epsilon\downarrow0$. \Cref{thm:gaud-volume-nef-limit}, with the Hermitian form $\kappa+w$, shows that the variational volume of $\theta_\pm(b)$ is one.
\end{proof}

\subsubsection{Current cones and endpoint equations}
\label{sec:warped-current-endpoint-equations}

Retain $X=B\times C$, $\theta=b\kappa+(ca+e)w$ with real parameters $b,c,e$, and the probability $\mu=\sigma\wedge w$ of \cref{sec:warped-gauduchon-completion}. Localized closed tests determine its current cone as well as its smooth cone. At a finite unit-volume end, the reciprocal coefficient is integrable with an exponent greater than one, so the base equation still has a continuous solution. Write $h=\kappa+w$.

\Needspace{14\baselineskip}
\begin{proposition}\label{prop:warped-current-cones}\index{Cone duality!warped product currents}
In this real span, the pseudoeffective and nef cones both equal
\[
 \{b\geq0,\ e\geq0,\ e+cM\geq0\}.
\]
A background admits a strict positive current representative if and only if $b>0$, $e>0$ and $e+cM>0$. Thus its strict-current cone equals the smooth-positive cone $\mathscr C$ in \cref{thm:warped-gauduchon-volume}.
\end{proposition}
\begin{proof}
Put $L=ca+e$. Suppose an arbitrary quasi-plurisubharmonic $u$ satisfies $\theta+\ddc u\geq\delta h$, where $\delta\geq0$. The closed semipositive test $\kappa^{m-1}\wedge w$ and distributional Stokes give $b\geq\delta$. For every smooth nonnegative top-degree form $\eta$ on $B$, its pullback to $X$ is closed by dimension, and
\[
 \int_X(\theta+\ddc u-\delta h)\wedge\eta
   =\int_B(L-\delta)\eta\geq0.
\]
If $L<\delta$ at a point, a nonzero test supported in that open set contradicts this inequality. Hence $L\geq\delta$ everywhere. For $\delta=0$ these are precisely $b\geq0$, $e\geq0$, $e+cM\geq0$; conversely they make $\theta$ semipositive. Such forms are nef because $\theta+\varepsilon h>0$ for every $\varepsilon>0$. Conversely, applying the same tests to smooth positive representatives of $\theta+\varepsilon h$ and sending $\varepsilon$ to zero gives those inequalities. For $\delta>0$ they force the three strict inequalities, and the converse follows because $\theta$ itself is then strictly positive. No restriction of $u$ to an extreme fiber is needed.
\end{proof}

\Needspace{14\baselineskip}
\begin{proposition}\label{prop:warped-continuous-endpoint-equation}\index{Monge--Ampere equation@Monge--Amp\`ere equation!finite warped boundary}\index{Regularity!continuous solution without C11@continuous solution without $C^{1,1}$}
Under the extreme-set hypotheses of \cref{lem:warped-endpoint-length}, let $\theta=b\kappa+L_0w$ be either finite face $\theta_+(b)$ or $\theta_-(b)$ of \cref{thm:warped-unit-volume-completion}, with $b>0$. There is a continuous $\theta$-plurisubharmonic potential $u$, pulled back from $B$, such that
\begin{equation}\label{eq:warped-continuous-endpoint-equation}
 (\theta+\ddc u)^n=\mu.
\end{equation}
The product is the ordinary bounded-potential product. No solution of this equation is globally $C^{1,1}$.
\end{proposition}
\begin{proof}
Let $f=a$ at the first face and $f=1-a/M$ at the second, and let $r>1$ be the corresponding codimension-to-order ratio. The coefficient $L_0$ is a positive constant multiple of $f$. In each extreme tube the integral of $f^{-p}$ has the same convergence condition as
\[
 \int_0^{s_0}s^{\operatorname{codim}_{\mathbb R}Z_i-1-k_ip}\,\mathrm{d}s,
\]
where $s$ is normal radius and $s_0>0$ is small. Thus $1/L_0\in L^p(B,\sigma)$ for any fixed $1<p<r$. The zero set has smooth volume zero, and the density can be assigned any finite value there. The face formulas give
\[
 H_0:=\int_B\frac{\sigma}{L_0}=nb^m.
\]
Apply \cref{thm:kn-continuous} on the positive closed background $b\kappa$ to obtain a continuous $\phi$ with $A=b\kappa+\ddc\phi\geq0$ and $A^m=\lambda\sigma/L_0$ for a scalar $\lambda>0$. The bounded closed mass identity in \cref{thm:mass-invariance-characterization} gives $\int_BA^m=b^m$, hence $\lambda=b^m/H_0$. Set $u=\phi\circ p_B$, where $p_B:X\to B$ is projection. Its solution current is $A+L_0w\geq0$.

The pullback identity of \cref{lem:bounded-products-submersion} gives $A^{m+1}=0$ on $X$, and $(L_0w)^2=0$. The multilinear expansion in \cref{prop:bounded-products} therefore gives
\[
 (A+L_0w)^n=nL_0 A^m\wedge w
       =\frac{nb^m}{H_0}\sigma\wedge w=\mu.
\]
The cancellation $L_0(\sigma/L_0)=\sigma$ holds as measures, since the zero set has smooth volume zero and $A^m$ has the displayed absolutely continuous density.

Suppose a solution $q$ were globally $C^{1,1}$. Since $\mu/h^n$ is smooth and strictly positive on compact $X$, \cref{prop:c11-ellipticity} would give $\theta+\ddc q\geq\delta h$ for some $\delta>0$. This contradicts \cref{prop:warped-current-cones}, because $L_0$ vanishes. The argument applies to every solution, not only to the pulled-back one.
\end{proof}

\subsubsection{Limits of the base potentials}
\label{sec:warped-endpoint-potential-limits}

The normalized potentials used to construct the interior balanced metrics have a unique continuous limit at each finite face. The compactness and uniqueness needed for this conclusion concern the closed base equation.

\Needspace{15\baselineskip}
\begin{proposition}\label{prop:warped-endpoint-potential-limits}\index{Gauduchon metric!boundary potential limits}\index{Compactness!warped endpoint potentials}
Under the hypotheses of \cref{thm:warped-unit-volume-completion}, let $\theta_j=b_j\kappa+L_jw\in\mathscr S$ converge in coefficients to a finite face $\theta_0=b\kappa+L_0w$. Normalize the smooth base potentials in \cref{thm:warped-gauduchon-volume} by $\sup_B\phi_j=0$ and put $u_j=\phi_j\circ p_B$. Then $u_j$ converge uniformly on $X$ to a continuous pulled-back solution $u_0$ of \eqref{eq:warped-continuous-endpoint-equation}. For every $1\leq k\leq n$,
\[
 (\theta_j+\ddc u_j)^k\rightharpoonup
             (\theta_0+\ddc u_0)^k.
\]
It is the unique normalized bounded solution pulled back from the base.
\end{proposition}
\begin{proof}
Write $v_j=c_j/e_j$ and $\psi_j=\phi_j/b_j$. The normalized closed base equation is
\[
 (\kappa+\ddc\psi_j)^m=F_{v_j}\sigma,
 \qquad F_v=\frac1{Q(v)(1+va)},\qquad\sup_B\psi_j=0.
\]
At the first face, $v_j\to\infty$. Set $\epsilon=1/v$ and $H_\epsilon=\int_B\sigma/(a+\epsilon)$. Then $F_v=[H_\epsilon(a+\epsilon)]^{-1}$ and $H_\epsilon\to H_a$. At the second face, $v_j\to-1/M$. With $f=1-a/M$ and $t=v+1/M>0$, the denominator is $Q(v)(f+ta)$ and $Q(v)\to H_-$. Thus the respective limits are
\[
 F_+=\frac1{H_a a},\qquad F_-=\frac1{H_-(1-a/M)}.
\]
Choose $1<p<r$, with $r>1$ the ratio for that face. Tubular integration as in \cref{prop:warped-continuous-endpoint-equation} gives $a^{-1}\in L^p(B,\sigma)$ or $f^{-1}\in L^p(B,\sigma)$, respectively. Near the first end, $F_v\leq C/a$; near the second, $F_v\leq C/f$. The extreme sets have smooth volume zero. Dominated convergence therefore gives $F_{v_j}\to F_\pm$ in $L^p(B,\sigma)$, and each limiting density has integral one. Since $\sigma/\kappa^m$ is smooth and strictly positive on compact $B$, the equation densities also have a common $L^p(B,\kappa^m)$ bound.

Apply \cref{thm:lp-uniform-compactness} to the fixed background $\kappa$. Every subsequence of $\psi_j$ has a uniformly convergent further subsequence. Bounded-product continuity identifies its limit with a normalized solution of $(\kappa+\ddc\psi_0)^m=F_\pm\sigma$. Closedness of $\kappa$ gives the Guan--Li condition, so \cref{uni-thm:GL-intro} makes this solution unique. Every cluster limit is consequently the same, which proves whole-sequence uniform convergence. Since $b_j\to b>0$, we obtain $u_j\to u_0=(b\psi_0)\circ p_B$ uniformly.

On each coordinate ball, coefficient convergence of the backgrounds gives a common smooth quadratic correction making all $u_j+q$ plurisubharmonic. These corrected potentials are uniformly bounded and converge uniformly. The capacity-continuity statement in \cref{sec:psh-local-capacity} gives weak convergence of all their bounded products, with uniform local mass bounds. Expanding off the correction as in \eqref{eq:bounded-expansion}, and using smooth convergence of the remaining background factors, gives the displayed current convergence. In top degree the interior products all equal $\mu$, so the limit solves the endpoint equation.

For any bounded pulled-back solution $v=\chi\circ p_B$ with $\sup_B\chi=0$, the horizontal current $A_\chi=b\kappa+\ddc\chi$ on $B$ is positive: pair its pulled-back current with positive base tests and $w$. The product expansion gives $nL_0A_\chi^m=\sigma$. Off the extreme set this determines $A_\chi^m=\sigma/(nL_0)$. That density already has total mass $H_0/n=b^m$, while closed bounded mass invariance gives $\int_BA_\chi^m=b^m$. Positivity excludes any residual mass on the extreme set. Thus $\chi/b$ satisfies the same normalized closed-base equation, and \cref{uni-thm:GL-intro} proves the asserted uniqueness.
\end{proof}

The statement selects the limit of the constructed base potentials. It does not identify arbitrary endpoint solutions or give a pointwise limit of the balanced metrics $\gamma_j=L_j^{-1/m}(b_j\kappa+\ddc\phi_j+L_jw)$.



\subsubsection{Traces of minimizing metrics}
\label{sec:warped-minimizer-trace-limits}

Uniform convergence of the base potentials does not control the minimizing balanced metrics themselves: by the cofactor formula \eqref{eq:warped-balanced-cofactor}, their $(n-1)$st powers contain the square of the reciprocal fiber coefficient. Depending on the ratio at the end, the traces of these powers stay bounded or blow up, and in the second case the trace-normalized powers concentrate on fibers over the extreme set.

Keep the sequence $\theta_j\to\theta_0$ of \cref{prop:warped-endpoint-potential-limits}, and let $\phi_0$ be the uniform limit of its base potentials $\phi_j$. Put $A_j=b_j\kappa+\ddc\phi_j$, and write $G_j=\gamma_j^m$, where $\gamma_j$ is the balanced minimizing metric of \cref{thm:warped-gauduchon-volume} for $\theta_j$. Traces are taken with respect to $h=\kappa+w$.

\Needspace{14\baselineskip}
\begin{proposition}\label{prop:warped-minimizer-trace}\index{Gauduchon metric!finite trace boundary}
Let $r>1$ be the ratio $r_0$ or $r_M$ of the end. After discarding finitely many terms, write $L_j=s_j(f+\epsilon_j)$ with $s_j>0$, $s_j\to s_0>0$ and $\epsilon_j\to0^+$. Here $f=a$, $s_j=c_j$ at the first face, and $f=M-a$, $s_j=-c_j$ at the second. Set $J_j=\int_B\sigma/L_j^2$. Then
\begin{equation}\label{eq:warped-minimizer-trace-mass}
 G_j=\frac{\sigma}{nL_j^2}+mA_j^{m-1}\wedge w,
 \qquad t_j:=\int_Xh\wedge G_j=\frac{J_j}{n}+mb_j^{m-1}.
\end{equation}
For $1<r<2$, $t_j\asymp\epsilon_j^{r-2}$; for $r=2$, $t_j\asymp\log(1/\epsilon_j)$. For $r>2$, the entire sequence converges weakly to the finite positive closed current
\[
 G_0=\frac{\sigma}{nL_0^2}+mA_0^{m-1}\wedge w,
 \qquad A_0=b\kappa+\ddc\phi_0,
\]
and $t_j\to t_0=\int_Xh\wedge G_0>0$.
\end{proposition}
\begin{proof}
On $\mathscr S$, $H(c_j,e_j)=nb_j^m$, so $A_j^m=\sigma/(nL_j)$. Substituting this in \eqref{eq:warped-balanced-cofactor} gives the formula for $G_j$. Stokes' theorem on $B$, for the closed form $\kappa$, gives $\int_B\kappa\wedge A_j^{m-1}=b_j^{m-1}$, which proves the trace identity. The factorization of $L_j$ uses $\epsilon_j=e_j/c_j$ at the first end and $\epsilon_j=e_j/(-c_j)-M$ at the second, and the coefficient limits give $s_j\to s_0>0$.

On each tube around an extreme component, $f\asymp\operatorname{dist}(Z_i)^{k_i}$, so the integral of $f^{-2}$ is comparable to
\[
 \int_0^{\varrho_0}\frac{\varrho^{\operatorname{codim}_{\mathbb R}Z_i-1}}{(\varrho^{k_i}+\epsilon_j)^2}\,\mathrm{d}\varrho,
\]
where $\varrho$ is normal radius, $\varrho_0>0$ is small and $r_i=\operatorname{codim}_{\mathbb R}Z_i/k_i$. Splitting the integral at $\varrho=\epsilon_j^{1/k_i}$ gives order $\epsilon_j^{r_i-2}$ for $r_i<2$, order $\log(1/\epsilon_j)$ for $r_i=2$, and a finite limit for $r_i>2$. The finite sum over components is governed by $r=\min_i r_i$. Since $b_j$ and $s_j$ have positive finite limits, these estimates prove the three alternatives for $t_j$.

If $r>2$, then $f^{-2}$ is integrable and dominates $L_j^{-2}$ up to a fixed constant factor, so $\sigma/L_j^2\to\sigma/L_0^2$ in total variation. Since $w\wedge w=0$, we have $A_j^{m-1}\wedge w=(\theta_j+\ddc u_j)^{m-1}\wedge w$, which converges weakly by \cref{prop:warped-endpoint-potential-limits}. This proves $G_j\rightharpoonup G_0$, and $G_0$ is positive and closed as a weak limit of positive closed forms. Pairing with $h$ gives the limit of the traces, which is positive because of the term $\sigma/(nL_0^2)$.
\end{proof}

\Needspace{15\baselineskip}
\begin{proposition}\label{prop:warped-normalized-test-limits}\index{Gauduchon metric!concentrating trace}\index{Gauduchon current!warped fiber limit}
Suppose $1<r\leq2$, and put $T_j=G_j/t_j$. Every weak cluster limit of $(T_j)$ has the form $p_B^*\nu$, where $\nu$ is a probability measure on $B$ supported on the extreme components with ratio $r$. Such a limit is positive and closed, has trace one relative to $h$, and satisfies $\theta_0\wedge p_B^*\nu=0$. If those components reduce to a single point $p$, the whole sequence converges to the current of integration $[\{p\}\times C]$.

The trace-normalized smooth metrics $\eta_j=t_j^{-1/m}\gamma_j$ satisfy
\[
 \int_X\theta_j\wedge\eta_j^m=D_\mu(\eta_j)
                  =\frac{nb_j^m}{t_j}\longrightarrow0.
\]
Both quantities have order $\epsilon_j^{2-r}$ when $1<r<2$, and order $1/\log(1/\epsilon_j)$ when $r=2$. Their quotient remains one.
\end{proposition}
\begin{proof}
Let $\nu_j=\sigma/(J_jL_j^2)$. These are probability measures on $B$, and
\[
 T_j=\frac{J_j}{nt_j}p_B^*\nu_j
                   +\frac m{t_j}A_j^{m-1}\wedge w.
\]
The first coefficient tends to one. The second current is positive with trace $mb_j^{m-1}/t_j\to0$, hence converges weakly to zero. Compactness of probability measures on $B$ and weak continuity of pullback under the submersion $p_B$ therefore show that every cluster limit is $p_B^*\nu$ for a cluster limit $\nu$ of $(\nu_j)$. Away from small tubes around the extreme set, the density $1/L_j^2$ is uniformly bounded, while $J_j\to\infty$; that region receives no mass in the limit. The tube integrals for components with ratio greater than $r$ are $o(J_j)$ by the preceding estimates, so only components with ratio $r$ keep mass.

For a limit $T=p_B^*\nu$, $\kappa\wedge T=0$ for degree reasons on $B$, while $L_0w\wedge T=0$ because $L_0$ vanishes on the support of $\nu$. Its trace is one because $\int_Cw=1$. A top-degree measure on $B$ is closed, so its pullback is closed as well. If the support is the single point $p$, every cluster limit of $(\nu_j)$ is $\delta_p$; compactness then gives convergence of the whole sequence, and $p_B^*\delta_p$ is the current of integration over the fiber $\{p\}\times C$.

By Stokes' theorem and \eqref{eq:warped-minimizer-trace-mass}, $\int_X\theta_j\wedge G_j=\int_X(A_j+L_jw)\wedge G_j=n\int_BA_j^m=nb_j^m$. Since $\gamma_j$ minimizes and $\lambda_\mu(\theta_j)=1$ on $\mathscr S$, also $D_\mu(\gamma_j)=nb_j^m$. Both are multiplied by $t_j^{-1}$ under the trace normalization, and the trace estimates give their limits and rates.
\end{proof}

In the projective strip of \cref{ex:warped-projective-strip}, both ratios are two and each extreme set is a single point. At each of the two ends, the normalized minimizing tests therefore converge to the current of integration over the corresponding fiber $\{p\}\times\mathbb P^1$, and the numerator and the denominator of the quotient both decay like $1/\log(1/\epsilon_j)$. This is convergence of currents; it does not give a pointwise limit of the metrics.


\subsubsection{Projective and torus examples}

\begin{example}\label{ex:warped-projective-strip}
On $B=\mathbb P^2$, take normalized Fubini--Study $\kappa$, $\sigma=\kappa^2$, $C=\mathbb P^1$, and
\[
 a([Z])=\frac{|Z_1|^2/2+|Z_2|^2}
                  {|Z_0|^2+|Z_1|^2+|Z_2|^2}.
\]
Then $M=1$, $r_0=r_M=2$, and the completion is a closed strip with finite faces
\[
 b\kappa+\frac{4\log2}{3b^2}a w,\qquad
 b\kappa+\frac{4\log2}{3b^2}(1-a)w.
\]
For $v\ne0$,
\[
 Q(v)=\frac4{v^2}\bigl[(v+1)\log(1+v)
                      -(v+2)\log(1+v/2)\bigr],
 \qquad Q(0)=1.
\]
The width of the strip is the finite positive integral
\[
 \int_{-1}^{\infty}
       \sqrt{\frac{Q''/Q-2(Q'/Q)^2}{3}}\,\mathrm{d}v.
\]
\end{example}
\begin{proof}
The extreme sets are the points $[1:0:0]$ and $[0:0:1]$. In affine coordinates, $a$ and $1-a$ vanish there to exactly second order, so both real codimensions are four and both orders are two. To compute the law of $a$ under $\sigma$, polar integration in the chart $Z_0=1$ gives $2\,\mathrm{d}s_1\mathrm{d}s_2/(1+s_1+s_2)^3$, where $s_j=|z_j|^2$. The transformation $u_j=s_j/(1+s_1+s_2)$ gives the uniform density two on the triangle $u_1,u_2\geq0$, $u_1+u_2\leq1$, and $a=u_1/2+u_2$. For $0\leq s\leq1/2$, the triangle $u_1/2+u_2\leq s$ has area $s^2$, so the law of $a$ has density $4s$ on $[0,1/2]$. With $u_0=1-u_1-u_2$, exchanging $u_0$ and $u_2$ preserves the uniform measure and replaces $a$ by $1-a$; hence the law has density $4(1-s)$ on $[1/2,1]$. Integrating against these two densities gives $H_a=H_-=4\log2$ and the formula for $Q$, and the integral definition of $Q$ shows that $Q(0)=1$ is its continuous extension. \Cref{thm:warped-unit-volume-completion} gives the boundary lines and the finite width, with Euclidean coordinates $Y=\sqrt2\log b$ and $R(v)$.
\end{proof}

\begin{example}\label{ex:warped-torus-plane}
Let $E=\mathbb C/(\mathbb Z+i\mathbb Z)$, $B=E^m$, and write $z_j=u_j+i\tau_j$. Put
\[
 \eta_j=\frac i2\mathrm{d}z_j\wedge\mathrm{d}\bar z_j,
 \qquad s=(m!)^{-1/m},\qquad
 \kappa=s\sum_j\eta_j,\qquad\sigma=\kappa^m.
\]
For $a=\sin^2(\pi u_1)$, the unit-volume hypersurface is itself a complete Euclidean plane, with Euclidean coordinates
\[
 Y=\sqrt m\log b,\qquad
 R=\frac{\log(1+v)}{2\sqrt n},\qquad v>-1.
\]
In this example the minimizing balanced metrics of \cref{thm:warped-gauduchon-volume} can be written down by integrating a periodic function twice.
\end{example}
\begin{proof}
The forms descend to $B$, and $\sigma=\eta_1\wedge\cdots\wedge\eta_m$ is a probability form. Both extreme sets are real hypersurfaces along which $a$, respectively $1-a$, vanishes to second order, so $r_0=r_M=1/2$. Splitting the period and substituting $t=\tan(\pi u)$ gives
\begin{align*}
 Q(v)&=\int_0^1\frac{\mathrm{d}u}{1+v\sin^2(\pi u)}\\
     &=\frac2\pi\int_0^\infty\frac{\mathrm{d}t}{1+(1+v)t^2}
      =\frac1{\sqrt{1+v}}.
\end{align*}
Let $F_v(u)=[Q(v)(1+v\sin^2(\pi u))]^{-1}$, which has mean one over a period. Integrating the equation
\[
 \phi_v''=4\pi bs(F_v-1)
\]
twice, whose right side has mean zero, gives a smooth periodic potential $\phi_v$ of $u=u_1$; the first integration constant is chosen so that $\phi_v'$ has mean zero. In our normalization $\ddc\phi_v=\phi_v''\eta_1/(4\pi)$, hence
\[
 A=bs\left(F_v\eta_1+\sum_{j=2}^m\eta_j\right)>0,
 \qquad A^m=b^mF_v\sigma=\frac{b^m}{H}\frac{\sigma}{L}.
\]
Thus the minimizer of \cref{thm:warped-gauduchon-volume}, with its cofactor formula \eqref{eq:warped-balanced-cofactor}, is explicit in this example, without an appeal to \cref{thm:tw-smooth}. Differentiating $Q$ gives
\[
 \rho(v)=\frac1{4n(1+v)^2},\qquad
 g_{\mathscr S}=m\,\mathrm{d}U^2+
             \frac1{4n}\left(\frac{\mathrm{d}v}{1+v}\right)^2.
\]
As $U$ and $v$ range over $\mathbb R$ and $(-1,\infty)$, the coordinates $(Y,R)$ range over all of $\mathbb R^2$. Both boundary faces of the coefficient cone are therefore at infinite distance on the unit-volume hypersurface.
\end{proof}

These distances are those of the fixed three-parameter cone with the product probability $\mu=\sigma\wedge w$. The endpoint criterion of \cref{lem:warped-endpoint-length} uses the hypothesis that the extreme sets are disjoint unions of submanifolds along which $a$ vanishes to finite power order; intersecting extreme strata or vanishing of infinite order would require further estimates. The variances are computed from the minimizing metrics, but we do not assert that these metrics have a pointwise limit at a finite end.
